% Section 4.5, from lectures/L04/README.md: what you should be able to do, questions to test
% yourself, and the next lecture.

\section{Review}
\label{c4:sec:review}

\needspace{6\baselineskip}
\subsection*{What you should be able to do now}
\begin{itemize}
  \item Name the three pointer register pairs and say which registers each is made of.
  \item Choose between \code{ld}, \code{ld Z+}, \code{ld -Z} and \code{ldd Z+q} for a given job,
    and say what each leaves the pointer holding.
  \item Say why \code{ldd} exists, why it is limited to a displacement of 63, and why \code{X}
    cannot do it.
  \item Read a byte from program memory, and say why it takes a different instruction and a
    different kind of address.
  \item Lay out a structure by hand and reach any field of it in one instruction.
  \item Walk an array of structures, and say what the stride is and why it is not the field size.
  \item Say which addresses in the data space a structure may occupy, and which look usable and are
    not.
  \item Compute the deepest the stack will reach in a program with interrupts, and check that it
    does not meet your variables.
\end{itemize}

\needspace{6\baselineskip}
\subsection*{Questions to test yourself}
\begin{enumerate}
  \item \code{Z} is \code{r31:r30}. Which is the high byte, and what does that mean for
    \code{movw r30, r24}?
  \item \code{ld r16, Z+} and \code{ldd r16, Z+0} both read the byte \code{Z} points at. Give two
    differences.
  \item Why can \code{ldd} reach \code{Z+63} but not \code{Z+64}, and why is there no \code{ldd} for
    \code{X} at all?
  \item An array of LED structures starts at \code{0x0300}. Where does the fourth one start, and
    what is the arithmetic you would write to get there from the first?
  \item A great deal of AVR code allocates driver structures at \code{RAMEND + 1}. What is at that
    address on an ATmega328P, and what happens when you store there?
  \item Your program is two calls deep when an interrupt arrives, and the handler saves five
    registers and makes two more calls. How many bytes of stack are in use at the deepest moment?
\end{enumerate}

\needspace{6\baselineskip}
\subsection*{Looking ahead}
Chapter~5 is about counting time in hardware, and the arithmetic of not quite being able to:
\begin{itemize}
  \item Counting time in hardware instead of in a loop that does nothing else.
  \item Prescalers, and the five ratios you actually get.
  \item Turning a period you want into a register value that cannot represent it exactly, and
    saying by how much you missed.
  \item A blinker that keeps time while the processor does something else.
\end{itemize}
